\documentclass[conference]{IEEEtran}
\usepackage{amsmath,amssymb,amsfonts}
\usepackage{booktabs,tabularx,microtype}
\title{Problem Formulation for the Proactive Golden-Hour Bridge Model: UAV-Assisted Post-Earthquake Response}
\author{\IEEEauthorblockN{Working Thesis Draft}}
\begin{document}
\maketitle
\begin{abstract}
This paper formulates a proactive UAV-assisted supply-drop problem for post-earthquake response when ground rescue teams are limited or already committed. Scan UAV reports provide severity, confidence, and required item quantities for detected survivor-assistance locations, which are processed by a command center into candidate supply-drop tasks. The responder-UAV dispatch problem is modeled as a sequence of binary candidate-tour selection problems. Each candidate fixes an ordered task sequence, routed paths, completion times, candidate-specific loading, payload evolution, and energy requirements. The objective maximizes an assistance priority score evaluated at predicted service completion time. Persistent task reservations prevent duplicate dispatch across decision epochs, while measured resource updates separate planned consumption from actual battery and inventory changes. The resulting formulation provides a resource-aware planning model; its assistance scores are operational priorities rather than calibrated predictions of medical outcomes.
\end{abstract}

\section{Problem Statement and Assumptions}
After a major earthquake, scan UAVs survey the affected region and report possible survivor-assistance locations to a command center. For each detected location, the scan UAV provides three planning inputs: a severity score, a confidence score, and the required relief item types and quantities. The command center processes these reports by filtering low-confidence or duplicate detections and converting accepted detections into supply-drop tasks.

The response setting considered here is resource constrained. Ground rescue teams are limited in number and may already be committed to other high-priority locations, so immediate physical rescue cannot be assumed for every detected survivor. In this interval before ground teams can arrive, responder UAVs act as a proactive bridge by delivering essential supplies to selected detected locations. The objective is therefore not to replace ground rescue, but to decide which supply-drop tasks should be served by the available responder UAVs while ground rescue capacity is delayed or occupied.

At each ordinary decision epoch, the command center selects at most one tour per ready responder UAV and at most one new assignment per pending supply-drop task. Tours start and end at designated UAV bases, where selected candidates are loaded before departure. Active tours retain their reservations while ready UAVs are replanned, preventing the same task from being assigned repeatedly across planning cycles. The recourse extension later relaxes this commitment only for the unexecuted future of selected active tours and only after a task completion and accepted state update.

The baseline assumes stationary task locations, a static flight map within each planning cycle, centralized status reporting, and supplied demand and service-duration estimates. Confidence describes detection reliability; severity describes the relative importance of timely assistance. Each selected task is planned for full remaining supply delivery by one UAV. Unexpected partial or failed delivery is handled through execution feedback. Split-delivery optimization, arbitrary mid-flight diversion, and interrupted-service recourse are outside the first model. Simultaneous flight deconfliction is provided externally.

\section{Sets, Parameters, and System State}
Times use seconds on a common clock, energy uses joules, distance uses metres, and mass uses kilograms. Table~\ref{tab:notation} summarizes principal symbols.
\begin{table}[t]
\caption{Principal notation}\label{tab:notation}
\centering\footnotesize
\begin{tabularx}{\columnwidth}{@{}lX@{}}
\toprule
Symbol & Meaning\\\midrule
\multicolumn{2}{@{}l}{\textit{Sets and indices}}\\
$\mathcal I(t),\mathcal J,\mathcal Q$ & Known task IDs, UAV IDs, and item types\\
$\mathcal A(t),\mathcal T^{\mathrm{rec}}(t),\mathcal J_R(t)$ & Baseline eligible tasks, recourse task universe, and ready UAVs\\
$\mathcal J_A^{\mathrm{rec}}(t)$ & Active UAVs eligible for recourse\\
$i,j,q,r,\ell$ & Task, UAV, item, candidate, and leg indices\\
\midrule
\multicolumn{2}{@{}l}{\textit{Task parameters and state}}\\
$p_i$ & Task service/drop location\\
$c_i,\sigma_i$ & Confidence and severity scores in $[0,1]$\\
$r_i^q(t)$ & Remaining task demand\\
$\rho_i(t),o_i(t)$ & Task status and reservation owner\\
$g_i,d_i,D_i$ & Ground-arrival, latest useful, and enforced completion times\\
$\beta_i(s),H_i,\alpha_i,v_i(s)$ & Time-dependent benefit, half-life, decay rate, and priority score\\
\midrule
\multicolumn{2}{@{}l}{\textit{UAV parameters and state}}\\
$B_j(t),B_j^{\max}$ & Measured and maximum battery energy\\
$I_j^q(t)$ & Actual physical onboard inventory\\
$p_j(t),b_j$ & UAV position and UAV base\\
$K_j,W_j,\omega_q$ & Item-count capacity, payload-mass capacity, and item mass\\
$\eta_j(t),E_j^{\mathrm{res}}$ & UAV status and reserve energy\\
\midrule
\multicolumn{2}{@{}l}{\textit{Candidate quantities}}\\
$\mathcal C_j(t),\pi_{jr}$ & Feasible candidates and task sequence\\
$\mathcal C_j^R(t),\mathcal C_j^A(t)$ & Ready- and active-UAV recourse candidate sets\\
$\mathcal P(q_\ell,q_{\ell+1}),P^*_{jr\ell}$ & Feasible path set and selected shortest feasible leg path\\
$L(P),A(P),T_j(P)$ & Path length, positive ascent, and travel time\\
$R_{jr}^q,R_{jr}^{q,\mathrm{rec}}$ & Baseline candidate load and recourse future demand\\
$\widehat I_{jr\ell}^q,\widehat I_{jr\ell}^{q,A}$ & Ready- and active-UAV predicted onboard inventory\\
$\delta_{ij}(t),T_{jr},E_{jr},E_{jr}^{\mathrm{rem}}$ & Service duration, tour duration, tour energy, and remaining energy\\
$a_{ijr},V_{jr},V_{jr}^{\mathrm{rec}},\tau_{ijr}$ & Membership, baseline/recourse value, and completion time\\
\midrule
\multicolumn{2}{@{}l}{\textit{Decision and derived quantities}}\\
$\lambda_{jr},\mu_{jr},y_{ij},x_i$ & Baseline/recourse candidate-selection and derived assignment indicators\\
\bottomrule
\end{tabularx}
\end{table}

\subsection{Task State}
Task $i$ has detection time $t_i^{\mathrm{det}}$, accessible service location $p_i$, confidence $c_i$, severity $\sigma_i$, and remaining demands $r_i^q(t)\in\mathbb Z_{\geq0}$. Its status is
\[
\rho_i(t)\in\{\mathrm{pending},\mathrm{assigned},\mathrm{served},\mathrm{closed}\}.
\]
Its owner is $o_i(t)\in\mathcal J\cup\{\varnothing\}$. Closed tasks include verified false detections, ground-team takeovers, and expiration of a modeled service window; closure does not predict death. A pending or assigned supply-drop task is served only when all required item quantities have been delivered.

Given a ground-rescue arrival estimate $g_i$ and latest useful completion time $d_i$, define
\begin{equation}
D_i=\min\{g_i,d_i\}.
\end{equation}
Unavailable bounds are $+\infty$. If both are unknown, no finite service window is enforced and assistance before ground arrival is not guaranteed. Finite bounds are scenario assumptions, not universal medical golden-hour constants.

\subsection{UAV State}
UAV $j$ has position $p_j(t)$, base $b_j$, measured energy $0\leq B_j(t)\leq B_j^{\max}$, actual physical onboard inventory $I_j^q(t)\in\mathbb Z_{\geq0}$, empty mass $m_j^{\mathrm{empty}}$, capacities $K_j,W_j$, and reserve $E_j^{\mathrm{res}}>0$. Its status $\eta_j(t)$ is ready, active, turnaround, or unavailable. A ready UAV is at its base with resource measurements confirmed and is available for candidate-dependent loading. Charger and base-stock availability are external inputs; replenishment is not assumed unlimited.

\subsection{Global System State}
The state contains records, not just ID sets:
\begin{equation}
\begin{split}
X(t)=\big(&\{S_i(t):i\in\mathcal I(t)\},\{U_j(t):j\in\mathcal J\},\\
         &\mathcal V,\{\pi_j^{\mathrm{active}}(t):j\in\mathcal J\}\big).
\end{split}
\end{equation}
Records $S_i$ and $U_j$ contain all attributes, statuses, reservations, measurements, and active-tour progress defined above. Free space $\mathcal V$ accounts for UAV size and an obstacle-clearance margin. A detected position inside rubble must be mapped to an accessible drop/contact location, not treated automatically as a flyable point.

\subsection{Eligibility Sets}
The eligible-task and ready-UAV sets are
\begin{equation}
\begin{split}
\mathcal A(t)=\{i\in\mathcal I(t):\;&\rho_i(t)=\mathrm{pending},\ o_i(t)=\varnothing,\\
                               &c_i\geq c_{\min},\ t<D_i\},
\end{split}\label{eq:eligible}
\end{equation}
\begin{equation}
\mathcal J_R(t)=\{j\in\mathcal J:\eta_j(t)=\mathrm{ready},\ p_j(t)=b_j\}.
\end{equation}
Low-confidence tasks await verification. Repeated detections are deduplicated into existing records before insertion.

At a recourse epoch, assigned but unexecuted tasks must also be reconsiderable. Let $\mathcal J_A^{\mathrm{rec}}(t)$ denote the active UAVs whose most recent task completion has been accepted and whose current position, battery, inventory, and remaining mission state are confirmed. The recourse task universe is
\begin{equation}
\begin{split}
\mathcal T^{\mathrm{rec}}(t)=\mathcal A(t)\cup
\{i\in\mathcal I(t):\;&\rho_i(t)=\mathrm{assigned},\\
&o_i(t)\in\mathcal J_A^{\mathrm{rec}}(t),\\
&i\ \text{unserved}\}.
\end{split}
\end{equation}
Served tasks, closed tasks, physically completed tasks, and assigned tasks owned by active UAVs outside the recourse process are excluded. No displaced-task set is specified before optimization; retention, transfer, and release are derived from the selected recourse candidates.

\section{Time-Dependent Assistance Priority}
For each eligible supply-drop task, the benefit of UAV assistance is modeled as time dependent. For any completion time $s\geq t_i^{\mathrm{det}}$, define the remaining-assistance benefit
\begin{equation}
\beta_i(s)=\sigma_i e^{-\alpha_i(s-t_i^{\mathrm{det}})},\qquad
\alpha_i=\frac{\ln 2}{H_i}>0.
\end{equation}
The parameter $\sigma_i\in[0,1]$ represents the relative severity of the detected need, while $H_i$ is an assumed benefit half-life in seconds. Thus, $\beta_i(s)$ decreases exponentially with service delay and reaches half of its initial severity-weighted value after $H_i$ seconds. The quantity is used only as an operational benefit proxy; it does not model clinical deterioration directly.

Detection confidence is incorporated through the following assistance priority score:
\begin{equation}
v_i(s)=
\begin{cases}
\dfrac{2\beta_i(s)c_i}{\beta_i(s)+c_i},&\beta_i(s)+c_i>0,\\[4pt]
0,&\text{otherwise}.
\end{cases}\label{eq:value}
\end{equation}
This harmonic aggregation rewards tasks that are both time-critical and sufficiently reliable. A task receives a high score only when the remaining-assistance benefit and detection confidence are simultaneously high. The score satisfies $v_i(s)\in[0,1]$, but it is neither an expected number of lives saved nor a calibrated survival probability.

The score may be used to guide candidate generation, but it does not by itself determine visit order. In the optimization model, task value is evaluated at the predicted service completion time. Therefore, two tours visiting the same task set can receive different objective values if their completion times differ. Since confidence affects both eligibility and scoring, sensitivity analysis with respect to the confidence threshold is required during evaluation.

\section{Candidate Tour Construction}
\subsection{Candidate Tour Definition}
For $j\in\mathcal J_R(t)$, construct a finite feasible candidate set $\mathcal C_j(t)$. Candidate $r$ contains a nonempty ordered sequence
\begin{equation}
\pi_{jr}=(i_1,\ldots,i_{k_{jr}}),\qquad i_\ell\in\mathcal A(t),
\end{equation}
where $k_{jr}$ is the number of tasks in the candidate and all task IDs are distinct. Thus, $\pi_{jr}$ specifies both the selected tasks and their visit order. For example, $(i_1,i_2)$ and $(i_2,i_1)$ are different candidates because they can produce different travel times, completion times, and resource requirements.

Each tour starts from the designated base $b_j$, visits the tasks in the order given by $\pi_{jr}$, and finally returns to the same base. Tour nodes are $q_0=b_j$, $q_\ell=p_{i_\ell}$ for visits, and $q_{k_{jr}+1}=b_j$. Leg $\ell$ runs from $q_\ell$ to $q_{\ell+1}$, for $\ell=0,\ldots,k_{jr}$.

Equivalently, a candidate can be viewed as a precomputed bundle containing its task sequence, shortest feasible routed flight legs, completion times, required candidate load, resource requirements, and value. This baseline abstraction keeps the assignment model compact; later integrated formulations can replace the candidate bundle with explicit assignment, sequencing, and three-dimensional path decisions.

\subsection{Routed Paths and Travel Time}
For each leg from $q_\ell$ to $q_{\ell+1}$, let $\mathcal P(q_\ell,q_{\ell+1})$ denote the set of feasible collision-free routed paths whose endpoints are $q_\ell$ and $q_{\ell+1}$ and whose segments lie inside the modeled free space $\mathcal V$. The routed path used by candidate $r$ is the shortest feasible path in terms of three-dimensional path distance:
\begin{equation}
P^*_{jr\ell}\in
\arg\min_{P\in\mathcal P(q_\ell,q_{\ell+1})} L(P),
\qquad \ell=0,\ldots,k_{jr}.
\end{equation}
If multiple shortest feasible paths exist, any deterministic tie-breaking rule may be used. This is a shortest-distance routing assumption, not a minimum-energy routing optimization.

Each selected path $P^*_{jr\ell}=(w_0,\ldots,w_N)$ has the required endpoints and satisfies
\begin{equation}
[w_{n-1},w_n]\subseteq\mathcal V,\qquad n=1,\ldots,N.
\end{equation}
This checks routed segments, not the straight line between service locations. Define the total routed distance and total positive ascent as
\begin{align}
L(P)&=\sum_{n=1}^{N}\|w_n-w_{n-1}\|_2,\\
A(P)&=\sum_{n=1}^{N}
\max\{0,\operatorname{alt}(w_n)-\operatorname{alt}(w_{n-1})\}.
\end{align}
Here $L(P)$ measures the three-dimensional path length, while $A(P)$ records upward motion for later energy calculations. Assumed constant cruise speed $0<v_j^{\mathrm{cr}}\leq v_j^{\max}$ gives
\begin{equation}
T_j(P)=\frac{L(P)}{v_j^{\mathrm{cr}}}.
\end{equation}
This travel-time model is an approximation, not a consequence of the speed limit alone. The shortest-distance path is used to make candidate generation deterministic; energy feasibility is still checked afterward using the selected path, ascent, payload evolution, service energy, and reserve requirement.

\subsection{Task Completion Times}
Let $\delta_{ij}(t)>0$ be full remaining supply-drop service duration, including approach stabilization, hovering, delivery release, and any required confirmation time. Let $T_j^{\mathrm{out}},T_j^{\mathrm{in}}\geq0$ be launch/recovery overheads not already included in path times. Completion times satisfy
\begin{equation}
\tau_{i_1jr}=t+T_j^{\mathrm{out}}+T_j(P^*_{jr0})+\delta_{i_1j}(t),
\end{equation}
so the first task is completed after the current time, launch overhead, travel from the base, and its service duration. For each later task,
\begin{equation}
\tau_{i_\ell jr}=\tau_{i_{\ell-1}jr}+T_j(P^*_{jr,\ell-1})+\delta_{i_\ell j}(t),
\quad \ell\geq2.
\end{equation}
Thus, each subsequent completion time is obtained by adding travel from the previous task and the service duration at the next task. The total duration of candidate tour $r$, including return to base, is
\begin{equation}
\begin{split}
T_{jr}={}&T_j^{\mathrm{out}}+T_j^{\mathrm{in}}
+\sum_{\ell=0}^{k_{jr}}T_j(P^*_{jr\ell})\\
&+\sum_{\ell=1}^{k_{jr}}\delta_{i_\ell j}(t).
\end{split}
\end{equation}
Therefore, $\tau_{ijr}$ represents the completion time of a visited task, whereas $T_{jr}$ represents the complete mission duration from launch through return and recovery.

\section{Candidate Feasibility}
\subsection{Candidate Loading Feasibility}
For candidate $r$ of UAV $j$, the planned load of item type $q$ is the total remaining demand over all tasks visited by the candidate:
\begin{equation}
R_{jr}^q=\sum_{\ell=1}^{k_{jr}}r_{i_\ell}^q(t),
\qquad q\in\mathcal Q.
\end{equation}
Candidate feasibility is checked against the carrying limits of the UAV, not against arbitrary pre-existing onboard stock:
\begin{equation}
\sum_q R_{jr}^q\leq K_j,\qquad
\sum_q \omega_qR_{jr}^q\leq W_j.
\end{equation}
In the baseline, a selected candidate is loaded with exactly $R_{jr}^q$ units of each item type before departure. Base-stock availability is treated as an external input or operational precondition rather than a depot inventory optimization.

\subsection{Payload Evolution}
After the first $\ell$ planned deliveries in candidate $r$, the predicted onboard inventory is
\begin{equation}
\widehat I_{jr\ell}^q=R_{jr}^q-\sum_{h=1}^{\ell}r_{i_h}^q(t)\geq0,
\quad \ell=0,\ldots,k_{jr}.
\end{equation}
An empty sum is zero, so $\widehat I_{jr0}^q=R_{jr}^q$. Payload mass and total flying mass on outgoing leg $\ell$ are
\begin{equation}
m_{jr\ell}^{\mathrm{pay}}=\sum_q \omega_q\widehat I_{jr\ell}^q,\qquad
m_{jr\ell}=m_j^{\mathrm{empty}}+m_{jr\ell}^{\mathrm{pay}}.
\end{equation}
This representation captures the reduction in payload mass after planned deliveries. Because the baseline loads exactly the planned candidate demand, predicted onboard stock is zero after all planned deliveries; actual leftover stock can still occur during execution if service is partial or interrupted.

\subsection{Energy Model and Feasible Tours}
For a path $P$ flown with total mass $m$, use the approximate flight-energy function
\begin{equation}
E_j^{\mathrm{fly}}(P,m)
=e_j^{\mathrm{dist}}(m)L(P)+e_j^{\mathrm{asc}}(m)A(P).
\end{equation}
Both nonnegative coefficients have units J/m. The ascent term is an incremental correction beyond the distance-based component. These functions must be calibrated from UAV specifications or empirical flight data.

Let $P_j^{\mathrm{srv}}(m)\geq0$ be service power in watts and $E_{jr}^{\mathrm{over}}\geq0$ launch/recovery energy not already included in path energy. The total planned energy of candidate $r$ is
\begin{equation}
\begin{split}
E_{jr}={}&E_{jr}^{\mathrm{over}}
+\sum_{\ell=0}^{k_{jr}}E_j^{\mathrm{fly}}(P^*_{jr\ell},m_{jr\ell})\\
&+\sum_{\ell=1}^{k_{jr}}P_j^{\mathrm{srv}}(m_{jr,\ell-1})\delta_{i_\ell j}(t).
\end{split}
\end{equation}
Using pre-delivery mass throughout service is an explicit approximation. A candidate is retained only if it satisfies battery reserve and completion-deadline requirements:
\begin{equation}
E_{jr}+E_j^{\mathrm{res}}\leq B_j(t),\qquad
\tau_{ijr}\leq D_i\quad(i\in\pi_{jr}).
\end{equation}
Return energy is already included in $E_{jr}$; the reserve $E_j^{\mathrm{res}}$ is additional contingency energy. Consequently, every candidate passed to the assignment optimizer has already satisfied membership, shortest feasible routed reachability, candidate-loading, payload, time-window, and battery requirements. Because routing is shortest-distance based, these feasibility checks do not imply that each leg is energy-minimizing.

\section{Assignment Optimization}
\subsection{Candidate Value and Decision Variable}
For every feasible candidate, define fixed assignment and value coefficients:
\begin{equation}
a_{ijr}=\mathbf1\{i\in\pi_{jr}\},\qquad
V_{jr}=\sum_{i\in\pi_{jr}}v_i(\tau_{ijr}).
\end{equation}
The coefficient $a_{ijr}$ indicates whether task $i$ appears in candidate $r$ for UAV $j$. The coefficient $V_{jr}$ is the total completion-time value of that candidate. Completion times are defined only for tasks visited by the candidate.

The binary decision variable $\lambda_{jr}$ selects the whole candidate tour. This single variable couples task assignment, visit order, routed paths, completion times, and resource feasibility.

\subsection{Primary Assignment Model}
The baseline assignment problem is
\begin{align}
\max_{\lambda}\quad&
\sum_{j\in\mathcal J_R(t)}\sum_{r\in\mathcal C_j(t)}V_{jr}\lambda_{jr}
\label{eq:objective}\\
\text{subject to}\quad&
\sum_{r\in\mathcal C_j(t)}\lambda_{jr}\leq1,
\quad j\in\mathcal J_R(t),\label{eq:one-tour}\\
&\sum_{j\in\mathcal J_R(t)}\sum_{r\in\mathcal C_j(t)}a_{ijr}\lambda_{jr}\leq1,
\quad i\in\mathcal A(t),\label{eq:one-task}\\
&\lambda_{jr}\in\{0,1\},\quad
j\in\mathcal J_R(t),\ r\in\mathcal C_j(t).
\end{align}
Constraint~\eqref{eq:one-tour} allows each ready UAV to execute at most one candidate tour. Constraint~\eqref{eq:one-task} ensures that each eligible task receives at most one new assignment at the current epoch. Selecting no candidate means no dispatch, with zero planned time and energy. The all-zero decision is feasible if no tour is available.

\subsection{Reporting and State-Update Indicators}
For reporting and state-update purposes, define derived indicators
\begin{equation}
y_{ij}=\sum_{r\in\mathcal C_j(t)}a_{ijr}\lambda_{jr},\qquad
x_i=\sum_{j\in\mathcal J_R(t)}y_{ij}.
\end{equation}
Here $y_{ij}=1$ means that task $i$ is newly assigned to UAV $j$ by the selected candidate, while $x_i=1$ means that task $i$ receives a new assignment. Constraints~\eqref{eq:one-tour}--\eqref{eq:one-task} imply binary reporting indicators on these domains. These variables are epoch-level assignment summaries, not persistent service flags; task owners and statuses record commitments across decision epochs.

\subsection{Solution Interpretation}
With candidates fixed, this is a binary linear set-packing problem. An exact solve is optimal over the supplied candidate set after each candidate leg has been assigned its shortest feasible routed path, but it is not necessarily optimal over all possible physical tours or over energy-minimizing routes. Exhaustive order/path enumeration may be expensive. Shortlists, a visit cap, or column generation are implementation options. A candidate visit cap is a computational restriction, not payload capacity. Implementations should report candidate-generation rules, runtime, solver gap, and comparisons with exhaustive small instances where tractable.

\section{Rolling-Horizon Resource-Constrained Recourse and Reassignment}
\subsection{Recourse Scope and Trigger}
The baseline assignment may be extended by a rolling-horizon recourse step after new information arrives. The first recourse version is restricted to post-task-completion epochs: after UAV $j$ completes a task, the command center accepts the corresponding task-status, battery, position, and inventory update, and only then may revise the unexecuted future of eligible active UAVs. A newly detected task may arrive earlier, but the revised mission commitment is delayed until such a permitted recourse epoch. Arbitrary mid-flight diversion, partial-leg interruption, and interruption of an ongoing service action remain outside this first model and are handled by external safety policies.

All accepted execution before the recourse epoch is immutable. Completed path segments, elapsed service actions, delivered supplies, consumed energy, served tasks, and closed tasks are not re-optimized or counted again. Recourse starts from the current accepted system snapshot: $p_j(t)$, $B_j(t)$, $I_j^q(t)$, task statuses, remaining demands, and the current free-space model.

\subsection{Ready- and Active-UAV Recourse Candidates}
At a recourse epoch, ready and active UAVs draw tasks from the same recourse task universe $\mathcal T^{\mathrm{rec}}(t)$, but their physical feasibility rules differ. For a ready UAV $j\in\mathcal J_R(t)$, construct $\mathcal C_j^R(t)$ as baseline-style candidates with sequences $\pi_{jr}^{R}$ starting from $b_j$, visiting a nonempty ordered subset of $\mathcal T^{\mathrm{rec}}(t)$, and returning to $b_j$. These candidates use exact candidate-dependent loading, the capacity checks in Section~V-A, and the full launch-to-return energy structure in Section~V-C. Thus a ready UAV may accept either a newly eligible pending task or an unexecuted task transferred away from an active UAV.

For an active UAV $j\in\mathcal J_A^{\mathrm{rec}}(t)$, construct $\mathcal C_j^A(t)$ as future candidates beginning at the current recourse position $p_j(t)$ rather than at the base. Candidate $r\in\mathcal C_j^A(t)$ contains an ordered future sequence
\begin{equation}
\pi_{jr}^{A}=(i_1,\ldots,i_{k_{jr}^{A}}),\qquad
i_\ell\in\mathcal T^{\mathrm{rec}}(t),
\end{equation}
possibly empty only for a service-free safe-return candidate. Each nonempty active candidate visits its selected future tasks, recomputes shortest feasible routed legs between consecutive future nodes, recomputes completion times from the recourse time $t$, and finally returns to $b_j$ unless an external recovery policy applies. The incumbent remaining sequence of UAV $j$ is included in $\mathcal C_j^A(t)$ only if it remains feasible under the current accepted state and environment.

\subsection{Active-UAV Inventory and Payload Evolution}
Active UAVs cannot reload at a recourse epoch. For an active candidate $r$, let the required future demand of item type $q$ be
\begin{equation}
R_{jr}^{q,\mathrm{rec}}=\sum_{\ell=1}^{k_{jr}^{A}}r_{i_\ell}^q(t).
\end{equation}
The candidate is inventory-feasible only if
\begin{equation}
R_{jr}^{q,\mathrm{rec}}\leq I_j^q(t),\qquad q\in\mathcal Q.
\end{equation}
Unlike a ready-UAV candidate, the predicted onboard inventory of an active UAV is initialized by actual current inventory, not by the demand of the revised candidate:
\begin{equation}
\widehat I_{jr\ell}^{q,A}=I_j^q(t)-\sum_{h=1}^{\ell}r_{i_h}^q(t),
\qquad \ell=0,\ldots,k_{jr}^{A}.
\end{equation}
Thus $\widehat I_{jr0}^{q,A}=I_j^q(t)$, and unused items do not disappear merely because the revised future task set needs fewer supplies than the UAV currently carries. Payload mass for active recourse candidates is computed from $\widehat I_{jr\ell}^{q,A}$, so leftover inventory contributes to outbound, inter-task, and return-leg energy until it is delivered or returned to base. The zero-stock-after-service property applies only to ready-UAV exact candidate-demand loading, not to active recourse.

\subsection{Remaining Mission Energy and Value}
For active recourse candidates, define $E_{jr}^{\mathrm{rem}}$ as the planned energy of only the remaining mission from the current recourse state. It includes future travel, future service, return to base, recovery energy when applicable, and no previously consumed launch, flight, or service energy. Feasibility requires
\begin{equation}
E_{jr}^{\mathrm{rem}}+E_j^{\mathrm{res}}\leq B_j(t).
\end{equation}
This preserves the principle that planned energy is not subtracted again from measured battery.

Let $\pi_{jr}^{\mathrm{rec}}$ denote the task sequence of candidate $r$ in the appropriate ready or active family. The future value of a recourse candidate is
\begin{equation}
V_{jr}^{\mathrm{rec}}=\sum_{i\in\pi_{jr}^{\mathrm{rec}}}v_i(\tau_{ijr}^{\mathrm{rec}}),
\end{equation}
where $\tau_{ijr}^{\mathrm{rec}}$ is recomputed from the current recourse time and the candidate's revised shortest feasible future paths. Already served tasks contribute no value in the recourse objective. The original launch-time value $V_{jr}$ is not reused for active incumbent plans; the incumbent remaining plan, if feasible, receives a newly computed current-state value.

\subsection{Joint Recourse Selection}
Let $\mathcal J^{\mathrm{rec}}(t)=\mathcal J_R(t)\cup\mathcal J_A^{\mathrm{rec}}(t)$, with $\mathcal C_j^{\mathrm{rec}}(t)=\mathcal C_j^R(t)$ for ready UAVs and $\mathcal C_j^{\mathrm{rec}}(t)=\mathcal C_j^A(t)$ for active UAVs. Binary variable $\mu_{jr}$ selects one future candidate. The recourse optimization is
\begin{align}
\max_{\mu}\quad&
\sum_{j\in\mathcal J^{\mathrm{rec}}(t)}
\sum_{r\in\mathcal C_j^{\mathrm{rec}}(t)}
V_{jr}^{\mathrm{rec}}\mu_{jr}\\
\text{subject to}\quad&
\sum_{r\in\mathcal C_j^{\mathrm{rec}}(t)}\mu_{jr}\leq1,
\quad j\in\mathcal J^{\mathrm{rec}}(t),\\
&\sum_{j\in\mathcal J^{\mathrm{rec}}(t)}
\sum_{r\in\mathcal C_j^{\mathrm{rec}}(t)}
a_{ijr}^{\mathrm{rec}}\mu_{jr}\leq1,
\quad i\in\mathcal T^{\mathrm{rec}}(t),\\
&\mu_{jr}\in\{0,1\}.
\end{align}
Here $a_{ijr}^{\mathrm{rec}}=1$ if task $i$ appears in recourse candidate $r$ of UAV $j$. This joint set-packing model naturally determines retained assignments, new assignments, transfers, substitutions, resequencing, and tasks left temporarily unassigned. No pairwise ``replace $B$ by $N$'' variables are required. If no useful-service candidate is feasible for an active UAV, a feasible safe-return candidate may be selected; if even safe return is infeasible, external recovery handling is invoked.

\subsection{Ownership Semantics}
Current ownership is provisionally protected while recourse is optimized. At atomic commitment, a task that remains in a selected candidate of its current owner keeps that owner; a task selected by a different UAV transfers owner atomically; a pending task selected by any UAV becomes assigned to that UAV; and a previously owned task not selected by any recourse candidate has its owner cleared and returns to pending if it remains eligible. Expired, invalid, or ground-taken tasks follow the existing closure logic. Therefore a displaced task is not predefined before optimization, is not guaranteed immediate reassignment, and cannot disappear from the system state.

\section{Decision-Epoch State Updates}
\subsection{Decision-Epoch Workflow}
Let $t^-$ and $t^+$ denote states immediately before and after an event or update. At an ordinary baseline decision epoch, the command center first processes new measurements and operational events. It then rebuilds the eligible-task and ready-UAV sets, generates feasible candidate tours, solves the assignment model, and commits the selected reservations. Detections, verification changes, service outcomes, turnaround completion, and route or resource alerts may trigger new decision cycles.

Because the formulation is online, decisions are based on the latest accepted snapshot of the system state. If task status, UAV status, battery, or base-loading availability changes before dispatch is committed, the affected candidate is revalidated before the reservation is installed. A recourse epoch is more restrictive: it is committed only immediately after a task completion and the corresponding accepted state update, and it may revise only future unexecuted assignments.

\subsection{Detection, Loading, and Dispatch Commitment}
New deduplicated tasks enter pending with no owner. Confidence/deadline updates change eligibility. A selected candidate $r$ for UAV $j$ is committed for dispatch only after base loading is performed and confirmed. Candidate selection alone does not physically change onboard inventory. After loading confirmation, set
\begin{equation}
I_j^q(t^+)=R_{jr}^q,\qquad q\in\mathcal Q.
\end{equation}
Then install the mission commitment and mark the UAV active:
\begin{align}
\eta_j(t^+)&=\mathrm{active},&
\pi_j^{\mathrm{active}}(t^+)&=\pi_{jr},\\
\rho_i(t^+)&=\mathrm{assigned},&
o_i(t^+)&=j\quad(i\in\pi_{jr}).
\end{align}
Assigned tasks stay excluded by Eq.~\eqref{eq:eligible} at every later epoch until served or released. Loading does not consume battery energy in the flight model unless it is represented through the external turnaround or launch/recovery terms:
\begin{equation}
B_j(t^+)=B_j(t^-).
\end{equation}
Forecast $\widetilde B_j^{\mathrm{return}}=B_j(t^-)-E_{jr}$ is separate and never overwrites measured battery.

\subsection{Execution Feedback}
For interval $[t_a,t_b]$ without charging, let $\Delta\widehat E_j(t_a,t_b)$ be actual energy since the last accepted update:
\begin{equation}
B_j(t_b)=B_j(t_a)-\Delta\widehat E_j(t_a,t_b).
\end{equation}
A direct measurement may instead replace the estimate. The model uses only one accepted resource update for a given interval, so planned tour energy is not subtracted again from the measured battery. A negative computed state flags inconsistent telemetry or operational failure rather than being silently clipped. Reserve violations trigger safe-return handling.

Supply delivery by owner $j$ records integer delivered quantities
\begin{equation}
0\leq d_{ij}^q\leq\min\{r_i^q(t^-),I_j^q(t^-)\},
\end{equation}
and updates both quantities once:
\begin{align}
I_j^q(t^+)&=I_j^q(t^-)-d_{ij}^q,\\
r_i^q(t^+)&=r_i^q(t^-)-d_{ij}^q.
\end{align}
Mark a task served and clear its owner only when all remaining item demands are zero:
\begin{equation}
\rho_i(t^+)=\mathrm{served},\quad o_i(t^+)=\varnothing
\quad\text{if}\quad \sum_{q\in\mathcal Q}r_i^q(t^+)=0.
\end{equation}

After incomplete service, retain the reservation during a committed retry. Once abandonment is acknowledged, clear the owner and return the task to pending if its window is open; otherwise close it. Recompute remaining service duration for future dispatch. Partial attempts do not each earn full task reward. Realized reward is recorded once at full completion and is zero outside the modeled service window.

\subsection{Recourse Commitment}
At a permitted recourse epoch, current ownership remains provisional until the selected recourse candidates are committed. For every task in $\mathcal T^{\mathrm{rec}}(t)$, compare its previous owner with its selected recourse owner, if any. Retained tasks keep their owners, transferred tasks change owners atomically, newly selected pending tasks become assigned, and previously owned but unselected tasks are released to pending when they remain eligible. Tasks that expire, are invalidated, or are taken over by ground teams are closed under the existing rules. Active UAV future plans are replaced only by their selected recourse candidates; all already completed actions remain fixed.

\subsection{Cancellation and Return}
Route failure, changed deadlines, ground-team takeover, or mission failure invokes an external safe-return or recovery policy. Outstanding reservations are released only after cancellation acknowledgment, preventing an old and a new owner from acting on the same task concurrently. Completed tasks remain served. Ground takeovers, invalid detections, and expired tasks are closed; active owners must acknowledge cancellation before their reservations clear. In-flight safety handling is external to the ready-UAV optimization model.

\subsection{Turnaround Update}
On return set $p_j(t^+)=b_j$, $\eta_j(t^+)=\mathrm{turnaround}$, clear the completed active tour, and account only for energy since the last telemetry update. Do not subtract full-tour energy again. Leftover inventory is measured or handled at the base. Charging and future candidate-dependent loading take externally specified turnaround time. Confirmed post-turnaround measurements satisfy
\begin{align}
0\leq B_j^{\mathrm{new}}&\leq B_j^{\max},\\
\sum_q I_j^{q,\mathrm{new}}&\leq K_j,\qquad
\sum_q \omega_qI_j^{q,\mathrm{new}}\leq W_j.
\end{align}
After operational checks, install the confirmed battery and inventory measurements and mark the UAV ready for a later candidate-dependent loading step. Charger or supply shortages delay readiness or reduce feasible loading; stock and energy are not created automatically.

\section{Limitations and Validation}
\subsection{Model Limitations}
The formulation is intended as a planning model for responder-UAV supply delivery, not as a clinical outcome model. It does not calibrate survival probabilities, optimize scout verification, optimize depot replenishment, guarantee fairness, or represent weather, moving obstacles, and communication-link uncertainty probabilistically. The recourse extension does not create contingency inventory, reload active UAVs, or guarantee that every displaced task is immediately reassigned. Simultaneous UAV safety also requires an external deconfliction layer. Cruise-time, service-time, and energy approximations require parameter evidence before experimental claims are made. No universal half-life, deadline, or measured performance improvement is asserted.

\subsection{Validation Protocol}
Before experiments, the implementation should specify parameter sources, confidence interpretation, item-demand assumptions, and supply-drop service modalities. Evaluation should compare completion-time score, completed task count, waiting time, deadline misses, and energy consumption against feasible nearest-task and severity-priority baselines. Sensitivity studies should vary benefit half-lives, confidence thresholds, service durations, reserve requirements, and candidate caps. Alternative weighted value-time-energy objectives may be tested here, but they are not part of the core assignment formulation.

The simulation should also test operational edge cases, including multiple-epoch reservations, repeated detections, partial delivery, detours around blocked straight-line paths, empty candidate sets, and conservation of battery and inventory through return. Recourse-specific tests should include new high-priority tasks after task completion, rejected low-value changes, incumbent infeasibility, safe-return-only cases, active UAVs with leftover inventory, transferred tasks, unreassigned displaced tasks, deadline expiration during recourse, and competition among UAVs for the same task. Planned objective scores should be reported separately from realized outcomes. Solver validation and simulation validation remain necessary before drawing experimental conclusions.

\section{Conclusion}
This formulation represents post-earthquake responder-UAV dispatch as a repeated candidate-tour selection problem for supply-drop assistance. Complete candidate tours couple task membership, visit order, shortest feasible routed travel, completion-dependent reward, payload evolution, and energy feasibility through one binary decision family. Persistent reservations protect assignments across decision epochs, while actual energy and delivery measurements drive physical state changes. The rolling-horizon recourse extension allows the unexecuted future of eligible active tours to be reconsidered after task completion using only current physical inventory and measured battery. The candidate-tour abstraction provides a clean baseline and can later be replaced by explicit assignment, sequencing, and three-dimensional path variables in an integrated optimization model. The model is therefore implementable as a resource-aware command-center planning layer, provided that external safety policies and empirical parameter calibration are supplied before experimental deployment.
\end{document}
